\section{Conclusions and future work}
\label{sec:conclusions}

The delivered model is experiment 13 -- a ResNet18-U-Net trained on random $1536^2$
windows of the native $2048^2$ frame -- with its inference configuration, selected
on the validation set, frozen on the assignment's score: validation \textbf{Dice
0.6815}, the highest of the study, PQ 0.393 at the same point, training at
\textbf{3.04~GB} and inference at \textbf{2.50~GB}, inside the optional GPU budget
on both counts. The submitted training and test notebooks reproduce the
configuration and the submission end to end.

Three lessons generalise beyond this dataset. \textbf{Measure what the score
measures, where it measures it}: selecting the operating point at native
resolution, against the official metric definition, and on the objective actually
graded was worth more than most training changes -- all free. \textbf{Pick the
metric that sees the failure mode}: Dice passed pairwise annotator agreement at the
second experiment and then stayed nearly flat; every real improvement was visible
only in PQ's recognition term, and an evaluation restricted to a single metric
would have declared the project finished ten experiments early. \textbf{A falsified
hypothesis is not a wasted run}: over twenty ideas were measured and rejected, and
each rejection redirected the effort -- experiment 3's failure exposed the
resolution problem, the failure of every post-hoc mechanism showed the residual
error lives in the model's probability support, and the cross-selection protocol
separated real gains from selection noise at the $\pm 0.005$ level.

The recognition gap that remains -- recall of small, faint filaments against
incomplete annotations -- is, on the evidence gathered here, a training-data
problem more than an architecture problem. Two directions follow directly: a
cross-validated protocol in which every observation trains a model, so that the
operating point is selected on $6.7\times$ more annotator entries and gains below
the current noise floor become measurable; and pseudo-labelling of the filaments
the model finds and no annotator marked, which the organisers have confirmed do
exist in the data.
